\chapter{Numerical Analyis}
\section{Polynomial Interpolation}
\begin{definition}
	The \tib{Lagrange interpolation} problem consists of finding a polynomial $p \in \cP_n$ such that $p(x_i) = y_i$ for some $a \leq x_0 < \hdots < x_n \leq b$ and $y_0, \hdots, y_n$.
\end{definition}
A convenient solution to this problem would be finding a set of polynomials $L_i$ such that $L_i(x_j) = \delta_{ij}$, since those would let us solve any interpolation problem using the right linear combination. A first idea is
\begin{align*}
	L_i(x) := \prod_{i \neq j} (x - x_j)
\end{align*}
which gives us $L_i(x_j) = 0$ if $i \neq j$ and $L_i(x_i) \neq 0$. Now, we only need to normalize these polynomials such that $L_i(x_i) = 1$. This can be done by setting
\begin{align*}
	L_i(x) := \prod_{i \neq j} \frac{x - x_j}{x_i - x_j}
\end{align*}
\begin{importanttheorem}{}{}
	Given some $x_0 < x_1 \hdots x_n$, the polynomials
	\begin{align*}
		L_i(x) := \prod_{i \neq j} (x - x_j)
	\end{align*}
	fulfill $L_i(x_j) = \delta_{ij}$. They are known as \tib{Lagrange polynomials}.
\end{importanttheorem}
\begin{corollary}
	The polynomial 
	\begin{align*}
		p(x) = \sum_{i = 0}^n y_i \cdot L_i(x)
	\end{align*}
	is the unique solution of the Lagrange interpolation problem.
\end{corollary}
\begin{exercise}
	Show that $\sum_{i = 0}^n L_i(x) = 1$ for all $x \in [a,b]$.
\end{exercise}
\begin{theorem}
	Let $f \in C^{n+1}([a,b])$ such that $f(x_i) = y_i$. Let $p$ be the solution of the corresponding Lagrange polynomial interpolation problem. Then for every $x \in [a,b]$, there exists a $\xi \in [a,b]$ such that
	\begin{align*}
		f(x) - p(x) = \frac{1}{(n+1)!} \cdot f^{(n+1)}(\xi) \cdot \prod_{j = 0}^n (x - x_j).
	\end{align*}
\end{theorem}
\begin{corollary}
	\begin{align*}
		\norm{f - p}_{\sup} \leq \frac{\norm{f^{(n+1)}}_{\sup}}{(n+1)!} \cdot (b-a)^{n+1}
	\end{align*}
\end{corollary}
\begin{corollary}
	Increasing the number of nodes $x_i$ or decreasing $b-a$ leads to uniform convergence of the interpolation polynomials to $f$, as long as the derivatives of $f$ don't grow rapidly.
\end{corollary}
In practice, this leads to uniform convergence for most $f$. Given uniformly spaced nodes, one counterexample is given by \tib{Runge's function} $f(x) = \frac{1}{1 + 25x^2}$.

Evaluating the interpolation polynomial can be difficult and unstable in practice. We now focus our attention on deriving a scheme that allows a stable evaluation of $p(x)$ in $\cO(n^2)$ time.

\begin{definition}
	Given nodes $x_0, \hdots, x_n$ and values $y_0, \hdots, y_n$ and $j \in [0..n]$, we define $p_{i,j}$ as the Lagrange interpolation polynomial fulfilling $p_{i,j}(x_k) = y_k$ for all $k \in [i..i+j]$
\end{definition}
Note that $p_{i,0}(x) = y_i$ and $p_{0,n}(x) = p(x)$.
\begin{importanttheorem}{Neville Scheme}{}
	The polynomials $p_{i,j}$ are given by $p_{i,0}(x) = y_i$ and
	\begin{align*}
		p_{i,j}(x) = \frac{(x-x_i) \cdot p_{i+1,j-1}(x) - (x-x_{i+j})p_{i,j-1}(x)}{x_{i+j} - x_i}
	\end{align*}
\end{importanttheorem}
Recursive implementations of the Neville scheme are highly inefficient. A simple successive "forwards" evaluation scheme leads to a runtime in $\cO(n^2)$.
\begin{definition}
	Given $x_0, \hdots, x_{n-1}$, the \tib{Newton basis} $q_0, \hdots, q_n$ of $\cP_n$ is given by
	\begin{align*}
		q_j(x) = \prod_{k = 0}^{j-1} (x - x_k)
	\end{align*}
\end{definition}
\begin{importanttheorem}{Divided Differences}{}
	Given $p \in \cP_n$ and $x_0, \hdots, x_n \in [a,b]$, we can find the representation $p(x) = \sum_{j = 0}^n \lambda_j q_j(x)$ of $p$ in the corresponding newton basis by setting $\lambda_{i,0} = y_i$ and iterating
	\begin{align*}
		\lambda_{i,j} = \frac{\lambda_{i+1,j-1} - \lambda_{i, j-1}}{x_{i+j} - x_i}
	\end{align*}
\end{importanttheorem}
$p$ can then be evaluated easily using the Horner Scheme
\begin{align*}
	p(x) = \lambda_0 + (x- x_0) \cdot \lr(\lambda_1 + (x - x_1) \cdot (\lambda_2 + \hdots (\lambda_{n-1} + (x - x_{n-1}\lambda_n))))
\end{align*}
This approach has the advantage of allowing us to easily add more nodes later.
\newpar
Now, we shift our attention to the problem of lowering interpolation error. One way of doing this is by cleverly choosing our nodes $x_i$.
\begin{definition}
	The \tib{Chebyshev-Polynomials} on $t \in [-1,1]$ are given by
	\begin{align*}
		T_n(t) = \cos(n \cdot \arccos(t)).
	\end{align*}
	The roots of a Chebyshev polynomial are known as \tib{Chebyshev nodes}.
\end{definition}
\begin{lemma}
	We have $\abs{T_n(t)} \leq 1$ for all $t \in [-1,1]$.
\end{lemma}
\begin{lemma}
	The Chebyshev polynomials are given by the recurrence relation
	\begin{align*}
		T_0(t) &= 1\\
		T_1(t) &= t\\
		T_{n+1}(t) &= 2t T_n(t) - T_{n-1}(t)
	\end{align*}
\end{lemma}
\begin{corollary}
	We have
	\begin{align*}
		T_{n+1}(x) = 2^n \prod_{j = 0}^n (x - t_j)
	\end{align*}
\end{corollary}
\begin{lemma}
	The roots of $T_n$ are given by 
	\begin{align*}
		t_j = \cos\lr(\frac{\lr(j + \frac{1}{2})\pi}{n})
	\end{align*}
	for $j \in [0..n-1]$.
\end{lemma}
\begin{lemma}
	The extreme values of $T_n$ are given by
	\begin{align*}
		s_j = \cos\lr(\frac{j\pi}{n})
	\end{align*}
	for $j \in [0..n]$.
\end{lemma}
\begin{theorem}
	Among all possible choices of nodes for polynomial interpolation, Chebyshev nodes minimize the supremum norm of the interpolation polynomial. If Chebyshev nodes are used, the interpolation error becomes
	\begin{align*}
		\norm{f - p}_{\sup} \leq 2^{-n} \frac{\norm{f^{(n+1)}}_{\sup}}{(n+1)!}
	\end{align*}
\end{theorem}

\begin{definition}
	The \tib{Hermite interpolation problem} consists of finding a polynomial $p \in \cP_N$ such that $p^{(k_i)}(x_i) = y^{(k_i)}_i$, where $N + 1 = \sum_{i = 0}^n (k_i + 1)$.
\end{definition}
\begin{proposition}
	Hermite interpolation has a unique solution. Given only a single node, the solution is the corresponding taylor polynomial.
\end{proposition}
\section{Spline Interpolation}
\begin{importantdefinition}{}{}
	A function $s : [a,b] \to \bR$ is known as a \tib{Spline of degree $m$ and order $k$} if $s \in C^k([a,b])$ is a piecewise combination of polynomials of degree $m$ over some partition $I_n$ of $[a,b]$ into $n$ nodes. The space of $m,k$-splines on $P$ is written $\cS^{m,k}(I_n)$. If no order is given, it is generally assumed that $k = m-1$.
\end{importantdefinition}
\begin{proposition}
	For $k \geq m$, we have $\cP_m|_{[a,b]} = \cS^{m,k}(I_n)$.
\end{proposition}
\begin{theorem}
	Linear spline interpolation is solved by $s = \sum_{i = 0}^n y_i \phi_i$, where $\phi_i$ are the \tib{hat functions}, given by linearly interpolating $\phi_i(x_j) = \delta_j$
\end{theorem}
\begin{theorem}
	\begin{align*}
		\dim \cS^{m,m-1}(I_n) = n + m
	\end{align*}
\end{theorem}
\begin{definition}
	Given a cubic spline interpolation problems, we generally want to fulfill one of the following boundary conditions:
	\begin{enumerate}
		\item \tib{Natural boundary conditions}, given by
		\begin{align*}
			s''(a) = s''(b) = 0
		\end{align*}
		\item \tib{Complete boundary conditions} or \tib{Hermite boundary conditions}, given by
		\begin{align*}
			s'(a) = \alpha,\\
			s'(b) = \beta
		\end{align*}
		for some $\alpha, \beta \in \bR$.
		\item \tib{Periodic boundary conditions} $s'(a) = s'(b)$, $s''(a) = s''(b)$
	\end{enumerate}
\end{definition}
If a thin, elastic board has its position fixed at certain points, the shape that the board takes is exactly the solution of the corresponding cubic spline interpolation problem.
\begin{theorem}
	Given a spline $s$ and any other function $g$ interpolating the same points and fulfilling the same boundary condtions, the curvature of $s$ is a lower bound to the curvature of $g$.
\end{theorem}
\begin{theorem}
	Let $s_1$ be the linear spline interpolating $f \in C^2([a,b])$. Then
	\begin{align*}
		\norm{f - s_1}_{L^2} \leq \frac{h^2}{2}\norm{f''}_{L^2}
	\end{align*}
	with $h = \max\lr{x_i - x_{i-1} \mid i \in [1..n]}$.
\end{theorem}
\begin{theorem}
	Let $s_3$ be a cubic spline interpolating $f \in C^2([a,b])$. Then
	\begin{align*}
		\norm{f - s_3}_{L^2} \leq \frac{h^4}{4}\norm{f^{(4)}}_{L^2}
	\end{align*}
	with $h = \max\lr{x_i - x_{i-1} \mid i \in [1..n]}$.
\end{theorem}
\section{Discrete Fourier Transform}
\begin{importantdefinition}{Real Trigonometric Interpolation}{}
	Given an even number of evenly spaces nodes $x_0, \hdots, x_{n-1}$ with $n = 2m$ even, the \tib{real trigonometric ineterpolation problem} consists of finding $a_k, b_k \in \bR$, $k \in [1..m-1]$ and $a_0, a_m \in \bR$ such that
	\begin{align*}
		T(x) = \frac{a_0}{2} + \sum_{k = 1}^{m-1} \lr(a_k \cos(kx) + b_k \sin(kx)) + \frac{a_m}{2} \cos(mx)
	\end{align*}
	fulfills $T\lr(\frac{2\pi j}{n}) = y_j$ for all $j$.
\end{importantdefinition}
Trigonometric interpolation can be written much more compactly in the complex case:
\begin{importantdefinition}{Real Trigonometric Interpolation}{}
	Given an even number of evenly spaces nodes $x_0, \hdots, x_{n-1}$ with $n = 2m$ even, the \tib{real trigonometric ineterpolation problem} consists of finding $c_k \in \bC$, $k \in [1..n-1]$ such that
	\begin{align*}
		T(x) = \sum_{k = 0}^{n-1} c_k e^{ikx}
	\end{align*}
	fulfills $T\lr(\frac{2\pi j}{n}) = y_j$ for all $j$.
\end{importantdefinition}
\begin{theorem}
	A set of coefficients $c_k$ solves the complex trigonometric interpolation problem iff. the coefficients $a_k, b_k$ defined by
	\begin{align*}
		a_0 &= 2c_0\\
		a_k &= c_k + c_{2m-k}\\
		b_k &= i(c_k - c_{2m-k})\\
		a_m &= 2c_m
	\end{align*}
	solve the corresponding real trigonometric interpolation problem.
\end{theorem}
If we write $y = (y_0, \hdots, y_{n-1})^\top$ and $\vec \omega_n^k = (\omega_n^{jk})_{j \in [1..n]}$, the complex trigonometric interpolation problem is given by the system of linear equations
\begin{align*}
	y = \sum_{k = 0}^{n-1} \beta_k \vec \omega_n^k
\end{align*}
\begin{importanttheorem}{}{}
	The family $(\vec \omega_n^k)_{k \in [1..n]}$ is an orthogonal basis of $\bC^n$ fulfilling
	\begin{align*}
		\scalar{\vec \omega_n^k}{\vec \omega_n^l} = n\delta_{kl}
	\end{align*}
	known as the \tib{Fourier basis}.
\end{importanttheorem}
Complex trigonometric interpolation thus comes down to performing a change of basis from the standard basis into the Fourier basis. We know from linear algebra that the matrix $T_n = (\vec \omega_n^0, \hdots, \vec \omega_n^n)$ is the basis change matrix from the Fourier basis into the standard basis. Thus, its inverse must be the basis change matrix from the standard basis into the Fourier basis.
\begin{exercise}
	Show $T_n^{-1} = \frac{1}{n} T_n^*$
\end{exercise} Thus, we get:
\begin{importanttheorem}{Discrete Fourier Transform}{}
	\begin{enumerate}
		\item Change of basis from the standard basis of $\bC^n$ into the Fourier basis, known as the \tib{discrete Fourier transform}, is given by
		\begin{align*}
			\beta = \frac{1}{n} T_n^* y
		\end{align*}
		\item Change of basis from the Fourier basis of $\bC^n$ into the standard basis, known as the \tib{inverse discrete Fourier transform} or as \tib{Fourier synthesis}, is given by
		\begin{align*}
			y = T_n \beta
		\end{align*}
	\end{enumerate}
\end{importanttheorem}
\section{Numerical Integration}
The goal of numerical integration, also known as \tib{quadrature}\sidenote{Based on the ancient greek practice of finding the area of a shape by constructing a square of equal area.} is numerical approximation of the integral operator
\begin{align*}
	I(f) = \int_a^b f(x) ~dx.
\end{align*}
\begin{definition}
	A \tib{quadrature formula} on an interval $[a,b]$ is a linear map $Q : C^0([a,b]) \to \bR$ of the form
	\begin{align*}
		Q(f) = \sum_{i = 0}^n w_i.f(x_i)
	\end{align*}
\end{definition}
\begin{definition}
	The number $$\norm{Q} = \frac{1}{b-a} \sum_{i = 0}^n \abs{w_i}$$ is known as the \tib{stability indicator} of $Q$.
\end{definition}
The Riemann integral is defined as the limit of the quadrature formula
\begin{align*}
	Q(f) = \sum_{k = 0}^{n-1} (x_{i+1} - x_i).f(x_i)
\end{align*}
as the distance between the nodes approaches zero.
\begin{definition}
	A quadrature formula $Q$ is called \tib{exact of degree $d$} if $Q(p) = I(p)$ for all $p \in \cP_d$.
\end{definition}
\begin{lemma}
	If a quadrature formula is exact of degree at least $0$, we have $\sum_{i = 0}^n w_i = b-a$.
\end{lemma}
\begin{lemma}
	Given $n = 2m$ equidistant nodes, a quadrature formula on these $n$ nodes that is exact of degree $2d$ is also exact of degree $2d+1$
\end{lemma}
\begin{importanttheorem}{}{}
	The \tib{Newton-Cotes-Formula}
	\begin{align*}
		Q(f) = \sum_{i = 0}^n \lr(\int_a^b L_i(x) ~dx).f(x_i)
	\end{align*}
	is exact of degree $n$. The case $n = 0$ is known as the \tib{middle point rule}, the case $n = 1$ is known as the \tib{trapeze rule}, and the case $n = 2$ is known as the \tib{Simpson rule}.
\end{importanttheorem}
\begin{theorem}
	The trapeze rule is given in closed form by
	\begin{align*}
		T(f) = \frac{b-a}{2}(f(a) + f(b))
	\end{align*}
\end{theorem}
\begin{theorem}
	The Simpson rule is given in closed form by
	\begin{align*}
		T(f) = \frac{b-a}{6}\lr(f(a) + 6f\lr(\frac{a+b}{2})+ f(b)).
	\end{align*}
	since we chose equidistant nodes $x_0, x_1, x_2$, it is exact of degree $3$.
\end{theorem}
\newpar
We can increase the accuracy of a quadrature formula by partitioning our interval $[a,b]$ into smaller intervals, and applying a quadrature formula to each one separately. We write this as
\begin{align*}
	Q^n(f) = \sum_{k = 1}^n Q_k(f|_{(x_{j-1},x_k)})
\end{align*}
A summed quadrature formula is called exact of degree $d$ if
\begin{align*}
	\abs{Q^n(f) - I(f)} = \cO\lr(\lr(\frac{b-a}{n})^d)
\end{align*}
\begin{proposition}
	The summed trapeze rule is exact of degree $2$. The summed Simpson rule is exact of degree $4$.
\end{proposition}
\begin{lemma}
	A quadrature formula with $n+1$ weights has exactness degree at most $2n + 1$.
\end{lemma}
\begin{lemma}
	A quadrature formula with $n+1$ weights is exact of degree $n$ if and only if
	\begin{align*}
		w_i = \int_a^b L_i(x) ~dx
	\end{align*}
	If the Quadrature formula is exact of Degree $2n$, we have $w_i > 0$ for all $i$.
\end{lemma}
If the error of a summed quadrature formula $T(h)$ is of order $h^\gamma$ for a $\gamma \in \bN$, we have
\begin{align*}
	T(h) = T(0) + \phi(h)
\end{align*}
with $\phi(0) = 0$ and $\abs{\phi(h)} \leq ch^\gamma$ for some $c \in \bR$. Taylor expansion of $\phi$ at $0$ gives
\begin{align*}
	\phi(h) = c_1h^\gamma + c_2 h^{\gamma + 1} + o(h^{\gamma + 1})
\end{align*}
for some $c_1, c_2 \in \bR$ and, i.e.
\begin{align*}
	T(h) = T(0) + c_1 h^\gamma + c_2 h^{\gamma + 1} + o(h^{\gamma + 1})
\end{align*}
If we evaluate at $h/2$, we get
\begin{align*}
	T(h/2) = T(0) + \frac{c_1}{2^\gamma} h^\gamma + \frac{c_2}{2^{\gamma + 1}} h^{\gamma + 1} + o\lr(\frac{h^{\gamma + 1}}{2^{\gamma + 1}})
\end{align*}
This lets us eliminate the $c_1 h^\gamma$-term by setting
\begin{align*}
	T^*(h) := \frac{T(h) - 2^\gamma T(h/2)}{1 - 2^\gamma} = T(0) + c_2 \frac{1 - 2^{-1}}{1 - 2^\gamma} \cdot h^{\gamma + 1} + o(h^{\gamma + 1}),
\end{align*}
which is a new summed quadrature formula with an error of $h^{\gamma + 1}$, instead of $h^\gamma$, which can be a significant improvement for $\gamma \ll 1$.
Given a non-polynomial function $f \in C^k([a,b])$ whose exact integral $I(f)$ is known, we can investigate the errors
\begin{align*}
	e_h = \abs{I(f) - Q^N(f)}
\end{align*}
for various $h > 0$. Since $e_h \approx c_1 h^\gamma$, we can evauate at two different step sizes $h, H > 0$ and get
\begin{align*}
	c_1 \approx \frac{e_h}{h^\gamma} \approx  \frac{e_H}{H^\gamma}
\end{align*}
and thus
\begin{align*}
	\gamma \approx \frac{\log(e_h/e_H)}{\log(h/H)} = \frac{\log(e_h) - \log e_H}{\log(h) - \log(H)}
\end{align*}
We can now try out various pairs $h,H$, and know that we have obtained the correct $\gamma \in h$ if these different pairs all give nearly identical answers.
\section{Nonlinear Optimization}
Let $U \subseteq \bR^n$. Then typical problems include:
\begin{enumerate}
	\item Find $x^* \in U$ such that $f(x^*) = 0$ for some $f : U \to \bR^m$.
	\item Find $x^* \in U$ such that $\displaystyle g(x^*) = \min_{x \in U} g(x)$ for some $g : U \to \bR$.
\end{enumerate}
In general, neither of these problems have closed form solutions. Instead, approximate solutions are found by iterative algorithms, i.e. we find $(x_k)_{k \in \bN}$ with $x_k \to x^*$.
\begin{definition}
	A numerical algorithm which returns a sequence $(x_k)_{k \in \bN}$ approximating a solution $x^*$ given some initial vector $x_0 \in U$ is known as:
	\begin{enumerate}
		\item \tib{Globally convergent}, if $x_k \to x^*$ for all $x_0 \in U$.
		\item \tib{Locally convergent}, if there exists some $\epsilon > 0$ such that $x_k \to x^*$ for all $x_0 \in U$ with $d(x_0, x^*) < \epsilon$.
	\end{enumerate}
\end{definition}
\begin{definition}
	If $x_k \neq x^*$ for all $k$, we can define the \tib{rate of convergence} of $x_k$ by saying that a locally convergent algorithm has \tib{convergence order $\alpha \geq 1$} if there exists a $q \in \bR$ such that the errors $\delta_k = d(x^* - x_k)$ fulfill
	\begin{align*}
		\limsup_{k \to \infty} \frac{\delta_{k+1}}{\delta_k^\alpha} = q.
	\end{align*}
	If $\alpha = 1$, we additionally demand $q \leq 1$.
\end{definition}
\begin{proposition}
	If $\Phi : \bR^n \to \bR^n$ is a contraction, the basic fixed point iteration $x_{k+1} = \Phi(x_k)$ is globally linearly convergent.
\end{proposition}
\begin{proposition}
	If $\Phi \in C^1(\bR, \bR)$ has a fixed point $x^*$, then the fixed point iteration $x_{k+1} = \Phi(x_k)$ is locally convergent if $\abs{\Phi'(x^*)} < 1$ and divergent if $\abs{\Phi'(x^*)} > 1$.
\end{proposition}
\section{Computing Roots Numerically}
\begin{theorem}
	\tib{Bisection Algorithm:}
	Let $f \in C^0([a,b], \bR)$ with $f(a)f(b) \leq 0$. Set $\epsilon > 0$. Consider the following algorithm:
	\begin{enumerate}
		\item Set $(a_0, b_0) = (a,b)$ and $k = 0$.
		\item Set $c_k = \frac{a_k + b_k}{2}$ and
		\begin{align*}
			(a_{k+1}, b_{k+1}) = \begin{cases}
				(a_k, c_k) & \tn{if} f(a_k)f(c_k) \leq 0\\
				(c_k, b_k) & \tn{otherwise}
			\end{cases}
		\end{align*}
		\item Return $b_{k+1}$ if $d(b_{k+1}, a_{k+1}) \leq \epsilon$. Otherwise set $k \gets k+1$ and return to step 2.
	\end{enumerate}	
	Then iteration stops after $J \leq 1 + \log_2\lr(\frac{b-1}{\epsilon})$ steps, the interval $[a_J, b_J]$ contains a root, and if $\epsilon = 0$, then the sequence $(c_k)_{k \in \bN}$ converges to a root linearly with $q = \frac{1}{2}$. 
\end{theorem}
\begin{theorem}
	\tib{Secant Method:} Let $f \in C^0([a,b] \to \bR)$. Then the following algorithm converges to a root of $f$, if one exists:
	\begin{enumerate}
		\item Let $\epsilon > 0$, $x_0 = a$, $x_1 = b$, and $k = 1$.
		\item If $f(x_k) \neq f(x_{k-1})$, set
		\begin{align*}
			x_{k+1} = x_k - \frac{x_k - x_{k-1}}{f(x_k) - f(x_{k-1})}f(x_k)
		\end{align*}
		\item Return $x_{k+1}$ if $d(x_{k+1}, x_k) \leq \epsilon$. Otherwise set $k \gets k+1$ and return to step 2.
	\end{enumerate}
\end{theorem}
\noindent
If $f \in C^1(U \subseteq \bR^n, \bR^n)$, Taylor approximation gives
\begin{align*}
	0 = f(x^*) = f(x) + Df(x)(x^*-x) + \phi(d(x^*, x))
	\approx f(x) + Df(x)(x^*-x),
\end{align*}
from which we can set $x := x_k$ to get a new iteration
\begin{align*}
	x_{k+1} = x_k - Df(x_k)^{-1} f(x_k) \approx x^*
\end{align*}
as long as $Df(x_k)$ is regular. This is the (multidimensional analogue of) the Newton method, and is generally much more efficient than the secant method.
\begin{theorem}
	Let $f \in C^2(U, \bR^n)$. Let $x^* \in U$ be a root of $f$ such that $Df(x^*)$ is invertible. Then there exists an $\epsilon > 0$ such that the Newton method converges for any $x_0 \in U$ with $d(x^*, x_0) < \epsilon$. Additionally, there exists a constant $c \geq 0$ such that the iterates $(x_k)_{k \in \bN}$ fulfill
	\begin{align*}
		\norm{x^* - x_{k+1}} \leq c \norm{x^* - x_k}^2
	\end{align*}
\end{theorem}
If $d(x^*, x_0) \geq \epsilon$, the Newton method might diverge. In this case, introducing a damping factor $\omega \in (0,1)$ and iterating
\begin{align*}
	x_{k+1} = x_k - \omega Df(x_k)^{-1}
\end{align*}
can lead to convergence, but makes convergence asymptotically slower (no longer quadratic in general).

